\section{实例一：RGB 彩色图像转灰度}

\subsection{数据并行分解}

红绿蓝 (red-green-blue, RGB) 图像的每个像素包含三个颜色分量。三个独立的一维数组 \texttt{red}、\texttt{green} 和 \texttt{blue} 分别保存三个输入分量，数组 \texttt{gray} 保存灰度输出。四个数组的长度均为 $WH$，其中 $W$ 为图像宽度，$H$ 为图像高度。

每个输出像素只依赖同一位置的三个颜色分量，不依赖其他像素。因此可把输出空间直接作为并行工作空间：
\[
\boxed{\text{一个线程} \longleftrightarrow \text{一个输出像素}}.
\]
线程在二维网格中先求自己负责的 $(r,c)$，再用 $p=rW+c$ 访问三个输入通道并写入一个输出值。

设某像素的红、绿、蓝分量分别为 $R_p,G_p,B_p$，灰度计算为
\[
\boxed{
Y_p=\frac{3}{10}R_p+\frac{6}{10}G_p+\frac{1}{10}B_p
}.
\]
三个系数之和为 1，因此当三个输入均位于 $[0,255]$ 时，浮点中间结果也位于该区间内。

\begin{examplebox}[title=单个像素的灰度计算]
若某像素满足
\[
R_p=100,\qquad G_p=150,\qquad B_p=200,
\]
则
\[
Y_p=0.3\times100+0.6\times150+0.1\times200
=30+90+20=140.
\]
该像素的三个输入值由一个线程读取，结果 140 由同一线程写入对应的灰度位置。
\end{examplebox}

\subsection{完整核函数}

核函数使用三个独立通道数组与上述权重。线性下标 \texttt{pix} 在边界判断通过后计算，使边界保护与内存访问的对应关系更直接。

\par\noindent\begin{minipage}{\linewidth}
\begin{lstlisting}[caption={带二维边界保护的 RGB 转灰度核函数}]
__global__
void rgb2gray_kernel(unsigned char* red,
                     unsigned char* green,
                     unsigned char* blue,
                     unsigned char* gray,
                     unsigned int width,
                     unsigned int height)
{
    unsigned int row = blockIdx.y * blockDim.y + threadIdx.y;
    unsigned int col = blockIdx.x * blockDim.x + threadIdx.x;

    if (row < height && col < width) {
        unsigned int pix = row * width + col;

        gray[pix] = red[pix]   * 3.0f / 10.0f
                  + green[pix] * 6.0f / 10.0f
                  + blue[pix]  * 1.0f / 10.0f;
    }
}
\end{lstlisting}
\end{minipage}\par

后缀 \texttt{f} 表示单精度浮点字面量。乘加表达式先产生浮点中间结果，再在赋给 \texttt{unsigned char} 时转换为字节值。源实现未加入额外舍入步骤，因此若结果带小数部分，赋值转换会去除其小数部分。

\subsection{启动配置}

采用 $16\times16$ 线程块时，二维启动配置可写为
\par\noindent\begin{minipage}{\linewidth}
\begin{lstlisting}[caption={RGB 转灰度核函数的二维启动配置}]
dim3 threadsPerBlock(16, 16);
dim3 numBlocks(
    (width  + threadsPerBlock.x - 1) / threadsPerBlock.x,
    (height + threadsPerBlock.y - 1) / threadsPerBlock.y);

rgb2gray_kernel<<<numBlocks, threadsPerBlock>>>(
    red, green, blue, gray, width, height);
\end{lstlisting}
\end{minipage}\par

该配置创建的矩形线程区域至少覆盖整个图像。右边界与下边界多出的线程仍会执行到条件判断，但不会读取或写入任何像素。

\subsection{一次线程索引的完整追踪}

考虑宽 $W=76$、高 $H=62$ 的图像，块尺寸为 $16\times16$。某线程满足
\[
\texttt{blockIdx}=(3,2),
\qquad
\texttt{threadIdx}=(5,7).
\]
其全局列、行坐标分别为
\[
\begin{aligned}
c&=3\times16+5=53,\\
r&=2\times16+7=39.
\end{aligned}
\]
因为 $53<76$ 且 $39<62$，该线程对应有效像素。线性下标为
\[
p=39\times76+53=3017.
\]
于是该线程读取
\[
\texttt{red[3017]},\quad
\texttt{green[3017]},\quad
\texttt{blue[3017]},
\]
并写入 \texttt{gray[3017]}。

若同一块中的线程改为 \texttt{threadIdx.y}=15，则行坐标变为
\[
r=2\times16+15=47,
\]
仍然有效；而位于最后一行块中的部分线程会得到 $r\geq62$，这些线程由外层条件排除。

\begin{keybox}
该实例包含一条完整的二维逐元素模式：
\[
\begin{aligned}
&\text{二维线程坐标}\\
&\quad\Downarrow\\
&\text{像素行列坐标 }(r,c)\\
&\quad\Downarrow\\
&\text{线性下标 }p=rW+c\\
&\quad\Downarrow\\
&\text{三个输入读取与一个输出写入}.
\end{aligned}
\]
所有四次内存访问共享同一个有效性条件 $r<H\land c<W$。
\end{keybox}
